\begin{filecontents*}{references.bib}
@article{carter1996,
  author = {Carter, Michael W. and Laporte, Gilbert and Lee, Sau Yan},
  title = {Examination Timetabling: Algorithmic Strategies and Applications},
  journal = {Journal of the Operational Research Society},
  year = {1996},
  volume = {47},
  number = {3},
  pages = {373--383},
  doi = {10.1057/jors.1996.37}
}

@article{burke2002,
  author = {Burke, Edmund K. and Petrovic, Sanja},
  title = {Recent Research Directions in Automated Timetabling},
  journal = {European Journal of Operational Research},
  year = {2002},
  volume = {140},
  number = {2},
  pages = {266--280},
  doi = {10.1016/S0377-2217(02)00069-3}
}

@article{qu2009,
  author = {Qu, Rong and Burke, Edmund K. and McCollum, Barry and Merlot, Liam T. G. and Lee, Sau Yan},
  title = {A Survey of Search Methodologies and Automated System Development for Examination Timetabling},
  journal = {Journal of Scheduling},
  year = {2009},
  volume = {12},
  number = {1},
  pages = {55--89},
  doi = {10.1007/s10951-008-0077-5}
}

@techreport{mccollum2007,
  author = {McCollum, Barry and McMullan, Paul and Burke, Edmund K. and Parkes, Andrew J. and Qu, Rong},
  title = {The Second International Timetabling Competition: Examination Timetabling Track},
  institution = {Queen's University Belfast},
  year = {2007},
  type = {Technical Report},
  number = {QUB/IEEE/Tech/ITC2007/Exam/v4.0/17},
  url = {https://www.eeecs.qub.ac.uk/itc2007/examtrack/report/Exam_Track_TechReportv4.0.pdf}
}

@article{mccollum2012,
  author = {McCollum, Barry and McMullan, Paul and Parkes, Andrew J. and Burke, Edmund K. and Qu, Rong},
  title = {A New Model for Automated Examination Timetabling},
  journal = {Annals of Operations Research},
  year = {2012},
  volume = {194},
  number = {1},
  pages = {291--315},
  doi = {10.1007/s10479-011-0997-x}
}

@article{leite2018,
  author = {Leite, Nuno and Fernandes, Carlos M. and Mel{\'\i}cio, Fernando and Rosa, Agostinho C.},
  title = {A Cellular Memetic Algorithm for the Examination Timetabling Problem},
  journal = {Computers \& Operations Research},
  year = {2018},
  volume = {94},
  pages = {118--138},
  doi = {10.1016/j.cor.2018.02.009}
}

@article{leite2019,
  author = {Leite, Nuno and Mel{\'\i}cio, Fernando and Rosa, Agostinho C.},
  title = {A Fast Simulated Annealing Algorithm for the Examination Timetabling Problem},
  journal = {Expert Systems with Applications},
  year = {2019},
  volume = {122},
  pages = {137--151},
  doi = {10.1016/j.eswa.2018.12.048}
}

@article{ceschia2023,
  author = {Ceschia, Sara and Di Gaspero, Luca and Schaerf, Andrea},
  title = {Educational Timetabling: Problems, Benchmarks, and State-of-the-Art Results},
  journal = {European Journal of Operational Research},
  year = {2023},
  volume = {308},
  number = {1},
  pages = {1--18},
  doi = {10.1016/j.ejor.2022.07.011}
}

@online{unitime2026,
  author = {{UniTime}},
  title = {Examination Timetabling},
  year = {2026},
  url = {https://www.unitime.org/uct_exams.php},
  urldate = {2026-05-05}
}
\end{filecontents*}

% ============================================================
%  ENSIA - Introduction to Artificial Intelligence
%  Project 1: Intelligent University Exam Scheduling System
%  Spring 2025-2026
% ============================================================

\documentclass[11pt, a4paper]{article}

% -- Packages ------------------------------------------------
\usepackage[T1]{fontenc}
\usepackage[utf8]{inputenc}
\usepackage{mathptmx}
\usepackage[top=2cm, bottom=1.5cm, left=1.5cm, right=1.5cm]{geometry}
\usepackage{setspace}
\usepackage{titlesec}
\usepackage{titletoc}
\usepackage{fancyhdr}
\usepackage{graphicx}
\usepackage{float}
\usepackage{booktabs}
\usepackage{array}
\usepackage{tabularx}
\usepackage{multirow}
\usepackage{amsmath, amssymb}
\usepackage{algorithm}
\usepackage{algpseudocode}
\usepackage{hyperref}
\usepackage{xcolor}
\usepackage{caption}
\usepackage{subcaption}
\usepackage{enumitem}
\usepackage{microtype}
\usepackage[backend=biber, style=apa, sorting=nyt]{biblatex}
\addbibresource{references.bib}

% -- Hyperref -----------------------------------------------
\hypersetup{colorlinks=true, linkcolor=black, citecolor=black, urlcolor=black}

% -- Section styling ----------------------------------------
\titleformat{\section}{\large\bfseries}{}{0em}{}[\titlerule]
\titleformat{\subsection}{\normalsize\bfseries}{\thesubsection.}{0.4em}{}
\titleformat{\subsubsection}{\normalsize\itshape}{}{0em}{}

% Tighter vertical spacing around headings to save page space
\titlespacing{\section}{0pt}{10pt}{4pt}
\titlespacing{\subsection}{0pt}{7pt}{3pt}
\titlespacing{\subsubsection}{0pt}{5pt}{2pt}

% -- Header / Footer ----------------------------------------
\pagestyle{fancy}
\fancyhf{}
\lhead{\small ENSIA -- Intro to AI}
\rhead{\small Intelligent Exam Scheduling System}
\cfoot{\small\thepage}
\renewcommand{\headrulewidth}{0.4pt}
\setlength{\headheight}{14pt}

% -- Spacing -------------------------------------------------
\setstretch{1.1}
\setlength{\parskip}{4pt}
\setlength{\parindent}{0pt}

% -- Caption -------------------------------------------------
\captionsetup{font=small, labelfont=bf, skip=4pt}

% ============================================================
\begin{document}

% ============================================================
%  COVER PAGE
% ============================================================
\begin{titlepage}
  \centering
  \vspace*{1cm}

  {\large\bfseries
    The National Higher School of Artificial Intelligence\\[0.3em]
    Introduction to AI Course Project\\[0.3em]
    Spring 2026
  }

  \vspace{1.8cm}
  \rule{\linewidth}{1pt}\\[0.5em]
  {\Large\bfseries Introduction to AI Project Paper}\\[0.5em]
  \rule{\linewidth}{1pt}

  \vspace{1.2cm}
  \noindent\textbf{\large Project Title:}\enspace Intelligent University Exam Scheduling System \underline{\hspace{9.6cm}}

  \vspace{1.2cm}
  \noindent\textbf{Student list:}

  \vspace{0.5em}
  \noindent
  \begin{tabular}{|p{9cm}|p{4.5cm}|}
    \hline
    \rule{0pt}{11pt}\textbf{Full name} & \textbf{Group}\\[2pt]
    \hline
    \rule{0pt}{16pt} Ilyes Oukil & G6 \\[5pt]\hline
    \rule{0pt}{16pt} Mohamed Boukra & G3 \\[5pt]\hline
    \rule{0pt}{16pt} Abdelkader Tarik Benabadji & G6 \\[5pt]\hline
    \rule{0pt}{16pt} Akram Benchouche & G3 \\[5pt]\hline
    \rule{0pt}{16pt} Abdelrahim Zemoura & G6 \\[5pt]\hline
    \rule{0pt}{16pt} Reda Amriou & G3 \\[5pt]\hline
  \end{tabular}
\end{titlepage}

% ============================================================
%  TABLE OF CONTENTS
% ============================================================
\newpage
\thispagestyle{fancy}
\tableofcontents

% ============================================================
%  ABSTRACT + INTRODUCTION
% ============================================================
\newpage
\setcounter{page}{1}

\phantomsection
\addcontentsline{toc}{section}{Abstract}

{\large\bfseries Abstract}
\vspace{2pt}
\titlerule
\vspace{6pt}

University examination timetabling is a recurrent and operationally
critical problem faced by higher education institutions at the end of
every academic semester. The task of assigning a set of exams to a
limited number of time slots and rooms, while simultaneously satisfying
a collection of hard and soft constraints, has been formally recognized
as a difficult combinatorial optimisation problem for which exhaustive
search becomes impractical at realistic scale. This report presents an
AI-driven examination scheduling system in which four complementary
problem-solving strategies are designed, implemented, and comparatively
evaluated. The current experimental study is built on the
\texttt{benchmark/1} instance distributed with the project repository,
derived from University of Nottingham semester-1 examination data for
1994--1995, and therefore reflects a real institutional scheduling
context rather than a purely synthetic toy example. Evaluation is
performed across multiple performance dimensions including hard
constraint violations, penalty cost, execution time, and room
utilisation. Results are expected to show clear trade-offs between
solution quality and computational effort, offering a practical basis
for choosing an algorithm according to institutional needs. An
interactive interface provides a Master Schedule grid view and a Student
Conflict Heatmap to support administrative decision-making.

\vspace{4pt}
\noindent\textbf{Keywords:} Combinatorial optimisation, Exam timetabling, A* search, Genetic Algorithm, Constraint Satisfaction Problem

\vspace{5pt}
\noindent\rule{\linewidth}{0.4pt}
\vspace{4pt}

\phantomsection
\addcontentsline{toc}{section}{Introduction (1 page max)}

{\large\bfseries Introduction}
\vspace{2pt}
\titlerule
\vspace{6pt}

The scheduling of university examinations is a real and recurring
administrative challenge faced by virtually every higher education
institution at the close of each academic semester. At its core, the
problem requires assigning each exam to a time slot and a room such that
no student is required to sit two exams simultaneously, no room exceeds
its seating capacity, and no two exams are placed in the same room at
the same time. While deceptively simple in statement, this assignment
task grows rapidly in complexity as student enrolments increase and
degree programmes become more modular, since both factors raise conflict
density, that is, the probability that two exams share at least one
common student~\parencite{carter1996,qu2009}.

From a computational standpoint, university examination timetabling is
closely related to graph coloring: exams can be represented as vertices
and a common enrollment induces an edge between two vertices, so
feasibility requires adjacent exams not to share the same time slot.
Because the underlying search space is combinatorial and heavily
constrained, no general polynomial-time exact method is known for the
realistic cases studied in practice~\parencite{carter1996,qu2009}. As a
result, the literature has progressively moved from simple constructive
heuristics toward richer search procedures, including constraint-based
methods, simulated annealing, tabu search, genetic algorithms, and
hybrid metaheuristics~\parencite{burke2002,qu2009,leite2019}.

Despite these advances, many institutions still depend on semi-manual
procedures, often starting from the previous semester timetable and
repairing conflicts by hand. Such workflows become fragile as cohorts
grow and room resources tighten. There is therefore a clear practical
need for automated scheduling systems that can rapidly construct valid
timetables, reduce administrative effort, and produce schedules that are
also more comfortable for students.

This project addresses that need by designing and implementing a
complete exam scheduling system around four distinct strategies: Greedy
search, Constraint Satisfaction (CSP), A* search, and a Genetic
Algorithm (GA). All four algorithms operate on a shared internal model
based on room--timeslot assignments and are compared under the same
evaluation function. For reproducible experimentation, the current
report uses the local \texttt{benchmark/1} dataset bundled with the
repository, a Nottingham-based real-world exam instance. Beyond
producing schedules, the system also provides an interactive interface
with a Master Schedule grid view and a Student Conflict Heatmap for
visual inspection of the timetable. The remainder of this report covers,
in order: related work, dataset construction, problem formulation and
solving techniques, application architecture, experimental results, and
finally discussion and conclusion.

% ============================================================
\newpage
\phantomsection
\addcontentsline{toc}{section}{State of the Art (1 page max)}
{\large\bfseries State of the Art}
\vspace{2pt}
\titlerule
\vspace{6pt}

Examination timetabling was first studied through constructive
heuristics inspired by graph coloring, because the central difficulty is
to separate mutually conflicting exams across a limited number of
periods. A landmark paper by \textcite{carter1996} compared several
algorithmic strategies on real institutional data and helped establish
shared benchmark culture in the field. Later surveys by
\textcite{burke2002} and \textcite{qu2009} show how the literature moved
from institution-specific heuristics toward reusable search procedures
and benchmark-driven empirical comparison.

Two broad families now dominate the literature. The first is
exact-or-model-based optimisation, such as integer programming,
constraint programming, and branch-and-bound. These approaches are
valuable when the scheduling model must express rich institutional rules
precisely, but they often struggle as problem size and constraint
interaction increase. This limitation motivated richer formal benchmark
models such as the ITC 2007 examination track and the later formulation
of \textcite{mccollum2012}, whose explicit goal was to reduce the gap
between simplified academic formulations and the constraints faced in
real deployments.

The second family is approximate search, especially metaheuristics and
hybrid methods. According to \textcite{qu2009}, tabu search, simulated
annealing, evolutionary algorithms, and memetic methods became popular
because they scale better on large instances and can improve a feasible
timetable without exhaustively exploring the search space. The ITC 2007
competition further accelerated this trend by standardising instances
and evaluation rules for realistic exam scheduling
problems~\parencite{mccollum2007}. More recent benchmark-oriented
surveys confirm that state-of-the-art results are still typically
achieved by highly engineered hybrid methods rather than by a single
simple heuristic alone~\parencite{ceschia2023}. Examples include the
FastSA simulated annealing variant of \textcite{leite2019} and the
cellular memetic algorithm of \textcite{leite2018}, both of which
report competitive results on established public benchmarks.

In parallel, production systems such as UniTime show what a full
institutional platform looks like: integrated exam data management,
distribution constraints, and solver support aimed at operational
deployment~\parencite{unitime2026}. Compared with such systems, our
project deliberately targets a lighter educational setting. Rather than
seeking the strongest published benchmark result, it compares four
transparent AI techniques under one common evaluator and couples them
with visual feedback tools, making the scheduling process easier to
understand, analyse, and discuss in a course-project context.

% ============================================================
\newpage
\phantomsection
\addcontentsline{toc}{section}{Data Set Building (1.5 pages max)}
{\large\bfseries Data Set Building}
\vspace{2pt}
\titlerule
\vspace{6pt}
\label{sec:data}

\subsection{Data Source and Collection}

The dataset used in this report comes from the local folder
\texttt{benchmark/1} included in the project repository. Its
\texttt{README} identifies the instance as the main exam timetabling data
for the University of Nottingham, semester 1, academic year 1994--1995.
The raw benchmark is therefore a real historical university dataset
rather than a synthetic classroom example. The same \texttt{README}
reports 800 exams, 7{,}896 students, 223 underlying courses, and
33{,}997 enrollment records.

The folder contains five complementary files: \texttt{README},
\texttt{students}, \texttt{exams}, \texttt{enrolements}, and
\texttt{data}. Student identifiers are already anonymised in the raw
files and appear only as alphanumeric codes, so no additional
de-identification step was needed for this project. To preserve
reproducibility, the benchmark was not manually altered; instead, the
backend parser converts the raw files into the internal JSON-like format
used by the scheduling algorithms and the frontend.

\subsection{Dataset Structure}

Each file contributes a different part of the scheduling instance. The
\texttt{exams} file is a fixed-width catalog containing exam codes,
titles, durations, and department identifiers. The \texttt{students}
file links each student to a course code, while the
\texttt{enrolements} file gives the student-to-exam relationships that
actually drive clash detection in the scheduler. Finally, the
\texttt{data} file stores the examination window, room capacities, and a
large set of additional benchmark-specific constraints such as room
preferences, coincidences, earliness priorities, ordering requirements,
and one room unavailability notice.

When parsed by the current backend, these files produce 800 active exams
with at least one enrolled student, 7{,}896 students, 16 rooms, and
33{,}997 enrollment links. Room capacities range from 15 to 270 seats.
The parser expands the calendar window from Monday 23 January 1995 to
Saturday 4 February 1995 into 32 candidate periods: three weekday
periods (09:00, 13:30, 16:30) and one Saturday morning period. This
period expansion follows the raw dates stored in the repository, even
though the Nottingham benchmark is sometimes presented in the literature
in a more compact period-based form.

\begin{table}[H]
  \centering
  \caption{Dataset Summary Statistics}
  \label{tab:dataset_stats}
  \begin{tabularx}{\linewidth}{lX}
    \toprule
    \textbf{Property}          & \textbf{Value} \\
    \midrule
    Number of courses / exams  & 800 exams \\
    Total students enrolled    & 7{,}896 students \\
    Number of rooms            & 16 rooms \\
    Number of time slots       & 32 candidate periods generated from the raw calendar window \\
    Conflict density           & 3.16\% conflicting exam pairs (10{,}113 conflicts over $\binom{800}{2}$ possible pairs) \\
    \bottomrule
  \end{tabularx}
\end{table}

\subsection{Preprocessing}

Preprocessing begins by parsing the fixed-width \texttt{exams} records
to recover each exam code, textual description, and duration. The
\texttt{enrolements} file is then scanned to build the set of students
registered in each exam and to derive the enrollment count used by the
solvers. From these student sets, a conflict graph is constructed: two
exams are adjacent if they share at least one student. In the current
instance this yields 10{,}113 conflicting exam pairs, corresponding to a
pairwise conflict density of approximately 0.0316.

The next step is resource generation. Room capacities are parsed from
the \texttt{data} file, and all room--timeslot pairs are created through
the Cartesian product of the 16 rooms and 32 generated periods, giving a
raw assignment domain of 512 possible placements per exam before
constraint pruning. This transformation is convenient for the four
implemented algorithms because they all manipulate a single integer
index that refers to one room--timeslot pair.

Finally, it is important to note what is \emph{not} yet encoded by the
current project. Although the raw \texttt{data} file contains fixed room
assignments, coincident exams, earliness priorities, ordering rules,
spread requirements, and a room unavailability constraint, the present
backend only uses the core information needed for room-capacity checks,
time generation, and conflict detection. This simplification makes the
benchmark easier to process inside a course project, but it also means
that \texttt{benchmark/1} acts as a demanding stress test rather than a
fully faithful implementation of every original Nottingham rule. In
particular, the simplified one-exam-per-room-slot model creates only 512
room-slot positions for 800 exams, so complete feasibility is not
guaranteed under this reduced formulation.

% ============================================================
\newpage
\phantomsection
\addcontentsline{toc}{section}{Problem Solving Techniques (2 pages max)}
{\large\bfseries Problem Solving Techniques}
\vspace{2pt}
\titlerule
\vspace{6pt}

\subsection{Problem Formulation}

Let $E = \{e_1, e_2, \dots, e_n\}$ be the set of exams, $R$ the set of
rooms, and $T$ the set of candidate periods. In the implementation, each
exam $e_i$ is represented by a decision variable $X_i$ whose value is
the index of one pair in the Cartesian product $R \times T$. For the
benchmark instance used here, the parser generates $|R| = 16$ rooms and
$|T| = 32$ periods, so each exam initially has up to 512 possible raw
placements before pruning. The student set attached to exam $e_i$ is
stored as $S_i$, which allows conflicts to be tested by set
intersection.

Three hard constraints are enforced by all four solvers:
\begin{enumerate}[leftmargin=1.2cm]
  \item No two exams may use the same room--timeslot pair.
  \item If two exams share at least one student, they may not be placed
  in the same timeslot.
  \item The assigned room capacity must be at least as large as the
  number of students registered in the exam.
\end{enumerate}

All algorithms are compared with the same evaluation function. For a
schedule $S$, the implemented fitness is
\[
F(S)=
\frac{\frac{1}{|A|}\sum_{i \in A} q_i}
{1 + \frac{L}{|A|} + 2\frac{P}{M}},
\qquad
q_i = -\frac{4n_i(n_i-c_i)}{c_i^2},
\]
where $A$ is the set of assigned exams, $n_i$ is the enrollment of exam
$i$, $c_i$ is the capacity of its assigned room, $L$ is the number of
late exams, $P$ is the total overlap count for pairs of exams sharing
students and scheduled on the same day, and $M$ is the total overlap
count over all exam pairs. If $M=0$, the last denominator term is
omitted in the code.

In simple terms, the evaluation function works as follows. It first
computes a room-allocation score for every assigned exam from the number
of enrolled students and the room capacity. It then divides that score
by two penalties: one for using late periods and another for placing
students with common registrations into multiple exams on the same day.
Higher fitness is therefore better. In the Genetic Algorithm, any hard
constraint violation immediately forces the fitness to $0$; in Greedy
and A*, partially assigned schedules are scored only on the exams that
have already been placed, while the CSP scores complete solutions found
during search.

\subsection{Greedy Search}

The Greedy solver uses a constructive one-pass strategy. In the current
implementation, exams are ordered by descending enrollment, so the
largest exams are scheduled first. Although the class comment mentions a
``Largest Enrollment + Degree'' heuristic, the actual sort key in code
is purely enrollment-based. This is a reasonable practical choice
because large exams are usually the hardest to fit into limited rooms.

For each exam, the algorithm scans every still-unused room--timeslot
pair and discards placements that violate a hard constraint. Among the
remaining candidates, it selects the one with minimum wasted seats,
while adding a very large local penalty of 1000 to late periods. The
method is therefore fast, deterministic, and easy to explain: it tries
to place the biggest exam first in the tightest feasible room and avoids
late slots whenever possible. Its weakness is myopia: once an exam is
placed, there is no repair step, so earlier choices can make later exams
hard or impossible to assign on dense instances.

\subsection{Constraint Satisfaction Problem (CSP)}

The CSP solver formulates timetabling as a backtracking search with one
variable per exam and one domain containing all currently feasible
room--timeslot indices. Search begins with the full domain for every
exam. At each step, the next variable is chosen by the Minimum
Remaining Values (MRV) heuristic, which selects the exam whose domain is
currently smallest. This focuses effort on the most constrained exams
first.

Once an exam is selected, candidate values are ordered by two practical
preferences implemented directly in the code: non-late periods are tried
before late ones, and rooms whose capacity is closer to the exam size
are tried before looser fits. After tentatively assigning a value, the
solver performs forward checking: every remaining exam has all
inconsistent values removed from its domain. If any future domain
becomes empty, the branch is abandoned immediately. The search is also
time-limited, so within the budget it keeps the best complete assignment
encountered according to the shared fitness function. This makes the CSP
component an informed depth-first search with strong pruning rather than
a plain exhaustive enumeration.

\subsection{A* Search}

The A* module performs search over \emph{partial} schedules. A state is a
tuple of length $n$ in which each position is either \texttt{None}
(unassigned) or the index of a room--timeslot pair. Before the search
starts, the solver computes a conflict degree for each exam, defined as
the number of other exams with which it shares at least one student. The
fixed expansion order then schedules exams by descending conflict degree
and, in case of ties, descending enrollment.

For a partial state at depth $d$, the heuristic is
\[
h(n) = (n-d)\times \overline{\deg}_{\text{conflict}} \times 0.1,
\]
that is, the number of remaining exams multiplied by the average
conflict degree in the instance. Candidate successors are filtered using
the same hard constraints as the other algorithms, ranked by wasted
seats plus a small late-slot cost, and only the best eight are expanded
in order to control branching. The priority queue also favors deeper
partial assignments, so in practice the method behaves as an A*-inspired
best-first search aimed at rapidly completing a feasible schedule rather
than as a textbook optimal A* with an admissible cost model. Complete
assignments are scored with the same fitness function used elsewhere.

\subsection{Genetic Algorithm (GA)}

In the Genetic Algorithm, one candidate timetable is encoded as a
chromosome whose length equals the number of exams. Each gene stores one
integer index into the room--timeslot list, so a full chromosome is a
complete tentative schedule. The initial population is generated
randomly over the full domain.

The implementation uses single-point crossover at the middle of the
chromosome: the first half is copied from one parent and the second half
from the other. Mutation is applied independently to each gene with a
user-provided probability; when mutation happens, the gene is simply
replaced by another random room--timeslot index. Parent selection is
based on roulette-wheel sampling, but one elite chromosome is preserved
unchanged in every generation.

Fitness evaluation is strict on feasibility. If two exams use the same
room--timeslot pair, if two conflicting exams share a timeslot, or if an
exam exceeds room capacity, the chromosome immediately receives fitness
$0$. Otherwise, the shared soft-constraint evaluator is applied. This
design makes the search easy to implement and keeps all algorithms on a
common comparison scale. The trade-off is that many random chromosomes
start infeasible, so population quality depends heavily on selection,
crossover, and repeated mutation over generations.

% ============================================================
\newpage
\phantomsection
\addcontentsline{toc}{section}{Application Design (2.5 pages max)}
{\large\bfseries Application Design}
\vspace{2pt}
\titlerule
\vspace{6pt}

\subsection{System Architecture}
\noindent\textit{[Architecture description goes here.]}

\subsection{Backend Modules}
\noindent\textit{[Backend module descriptions go here.]}

\subsection{Frontend \& Visualisation}
\noindent\textit{[Frontend description goes here.]}

\subsection{Technologies and Libraries}

\begin{table}[H]
  \centering
  \caption{Technologies Used}
  \label{tab:technologies}
  \begin{tabularx}{\linewidth}{llX}
    \toprule
    \textbf{Component}  & \textbf{Technology}    & \textbf{Purpose} \\
    \midrule
    Language            & Python 3.x             & Core implementation \\
    Backend             & FastAPI                & Scheduling API, CSV upload, and benchmark loading \\
    Notebook            & Jupyter                & Algorithm experiments \\
    Frontend            & HTML / CSS / JavaScript & Interactive controls and result display \\
    Visualisation       & Matplotlib / Seaborn   & Plots and heatmaps \\
    Data handling       & pandas                 & Dataset processing \\
    \bottomrule
  \end{tabularx}
\end{table}

% ============================================================
\newpage
\phantomsection
\addcontentsline{toc}{section}{Results and Analysis (2.5 pages max)}
{\large\bfseries Results and Analysis}
\vspace{2pt}
\titlerule
\vspace{6pt}

\subsection{Evaluation Metrics}
\noindent\textit{[Metric definitions go here.]}

\subsection{Experimental Setup}
\noindent\textit{[Experimental setup goes here.]}

\subsection{Comparative Results}

\begin{table}[H]
  \centering
  \caption{Performance Comparison of Scheduling Algorithms}
  \label{tab:results}
  \begin{tabularx}{\linewidth}{lXXXX}
    \toprule
    \textbf{Algorithm}    & \textbf{Hard Violations}
                          & \textbf{Penalty Cost}
                          & \textbf{Time (s)}
                          & \textbf{Utilisation (\%)} \\
    \midrule
    Greedy (Degree H.)    & -- & -- & -- & -- \\
    A*                    & -- & -- & -- & -- \\
    CSP + MRV + FC        & -- & -- & -- & -- \\
    Genetic Algorithm     & -- & -- & -- & -- \\
    \bottomrule
  \end{tabularx}
\end{table}

\noindent\textit{[Results discussion goes here.]}

\subsection{Key Insights}
\noindent\textit{[Key insights go here.]}

% ============================================================
\newpage
\phantomsection
\addcontentsline{toc}{section}{Discussion (1 page max)}
{\large\bfseries Discussion}
\vspace{2pt}
\titlerule
\vspace{6pt}

\noindent\textit{[Discussion text goes here.]}

% ============================================================
\vspace{10pt}
\phantomsection
\addcontentsline{toc}{section}{Conclusion (0.5 pages max)}
{\large\bfseries Conclusion}
\vspace{2pt}
\titlerule
\vspace{6pt}

\noindent\textit{[Conclusion text goes here.]}

% -- References ----------------------------------------------
\newpage
\printbibliography[heading=bibintoc, title={References}]

% ============================================================
%  APPENDICES
% ============================================================
\newpage
\appendix
\renewcommand{\thesection}{Appendix \Alph{section}}

\section{System Interface and Implementation}
\label{app:screenshots}

\noindent\textit{[Screenshots will be added before submission.]}

\newpage
\section{Individual Contributions}
\label{app:contributions}

\begin{table}[H]
  \centering
  \caption{Distribution of Responsibilities}
  \label{tab:contributions}
  \begin{tabularx}{\linewidth}{lX}
    \toprule
    \textbf{Full Name} & \textbf{Contributions} \\
    \midrule
    Ilyes Oukil  & Data collection; conflict matrix builder; Section~3. \\
    Mohamed Boukra  & Greedy and A* implementation; Sections 4.2 \& 4.4. \\
    Abdelkader Tarik & CSP implementation (backtracking + MRV + FC); Section 4.3. \\
    Akram Benchouche & Genetic Algorithm implementation; Section 4.5. \\
    Abdelrahim Zemoura  & Frontend \& visualisation; Section 5.3; Appendix A. \\
    Reda Amriou  & Team coordination; evaluation module; Sections 6 \& 7; References. \\
    \bottomrule
  \end{tabularx}
\end{table}

\end{document}
